% ============================================================
% ODE IVP Project — LaTeX Report Template
% Compile: pdflatex report.tex  (twice for the TOC)
% Structure follows the project brief: Background -> Model ->
% Methods -> Stability -> Numerical experiments -> Conclusions.
%
% This is the MERGED template: it combines the project-specific
% guidance for Project 1 (ODE IVPs) with the course-level report
% skeleton (algorithm environment, theorem environments, MMS
% convergence table, AI Transparency Log and ICS appendices).
% ============================================================

\documentclass[11pt,a4paper]{article}

\usepackage[utf8]{inputenc}
\usepackage[T1]{fontenc}
\usepackage{amsmath, amssymb, amsthm}
\usepackage{geometry}
\geometry{a4paper, margin=2.5cm}
\usepackage{graphicx}
\usepackage{booktabs}
\usepackage{listings}
\usepackage{xcolor}
\usepackage{hyperref}
\usepackage{enumitem}
\usepackage{float}
\usepackage{caption}
\usepackage{subcaption}
\usepackage{bm}
\usepackage{algorithm}
\usepackage{algpseudocode}
\usepackage{mdframed}

% --- colors ---
\definecolor{lecture}{RGB}{70,130,180}
\definecolor{seminar}{RGB}{205,92,92}
\definecolor{formbg}{RGB}{250,250,245}
\definecolor{formframe}{RGB}{180,170,140}

% --- code listing style (Python) ---
\lstset{
    language=Python,
    basicstyle=\ttfamily\scriptsize,
    keywordstyle=\color{blue}\bfseries,
    commentstyle=\color{green!40!black},
    stringstyle=\color{red!70!black},
    numbers=left,
    numberstyle=\tiny\color{gray},
    frame=single,
    breaklines=true,
    showstringspaces=false,
    tabsize=4,
    captionpos=b,
}

% --- theorem-like environments (course-level set) ---
\newtheorem{theorem}{Theorem}[section]
\newtheorem{lemma}[theorem]{Lemma}
\newtheorem{proposition}[theorem]{Proposition}
\newtheorem{corollary}[theorem]{Corollary}
\theoremstyle{definition}
\newtheorem{definition}[theorem]{Definition}
\newtheorem{assumption}[theorem]{Assumption}
\newtheorem{example}[theorem]{Example}
\theoremstyle{remark}
\newtheorem{remark}[theorem]{Remark}

% --- algorithm style ---
\algrenewcommand{\algorithmicrequire}{\textbf{Input:}}
\algrenewcommand{\algorithmicensure}{\textbf{Output:}}
\algrenewcommand{\algorithmiccomment}[1]{\hfill\textcolor{gray}{\# #1}}

% --- form box (for appendices) ---
\newmdenv[
    backgroundcolor=formbg,
    linecolor=formframe,
    linewidth=0.8pt,
    roundcorner=3pt,
    innermargin=5pt,
    outermargin=5pt,
    innertopmargin=8pt,
    innerbottommargin=8pt
]{formbox}

\title{\textbf{Numerical Solution of \texttt{<YOUR PROBLEM>}}\\[2mm]
\large Project 1 (ODE IVP) Report}
\author{
    Team \#\texttt{<N>} \\
    \small \texttt{<name1>, <id1> (Project Manager)} \quad
    \texttt{<name2>, <id2> (Mathematical Theory)} \\
    \small \texttt{<name3>, <id3> (Algorithm Implementation)} \quad
    \texttt{<name4>, <id4> (Visualization \& Report)} \\
    \small \texttt{<name5>, <id5> (Testing \& Validation)} \\
    \small Department of Applied Mathematics
}
\date{\today}

\begin{document}
\maketitle

\begin{abstract}
\noindent
[150--250 words answering: (1) what problem and why interesting; (2) what methods
you implemented; (3) key findings (convergence rates, stability observations);
(4) the main limitation or open question.]
\end{abstract}

\noindent\textbf{Keywords:} [e.g., Runge--Kutta methods, stiffness, von Neumann
stability, absolute stability, adaptive stepping]

\tableofcontents
\newpage

% ============================================================
\section{Introduction and Motivation}
% ============================================================

\subsection{Problem Background}
[Describe the physical, biological, or financial context of your chosen ODE
system. Why is this problem worth solving? Cite 2--3 references.]

\subsection{Mathematical Model}
[State the ODE system, initial conditions, and all parameter values. If the
system is stiff, say so and give the stiffness ratio.]

\begin{equation}
    \frac{dy}{dt} = f(t, y), \qquad y(t_0) = y_0.
    \label{eq:ivp}
\end{equation}

For the Robertson problem (example):
\begin{equation}
\begin{aligned}
    \dot{y}_1 &= -0.04\, y_1 + 10^4\, y_2 y_3, \\
    \dot{y}_2 &= \;\;0.04\, y_1 - 10^4\, y_2 y_3 - 3\times 10^7\, y_2^2, \\
    \dot{y}_3 &= \;\;3\times 10^7\, y_2^2,
\end{aligned}
\qquad y(0) = (1, 0, 0).
\label{eq:robertson}
\end{equation}

\subsection{Overview of the Report}
[A roadmap: ``In Section 2 we derive the discrete scheme and analyse stability.
In Section 3 we describe the implementation. Section 4 presents numerical
experiments. Section 5 concludes.'']

\newpage
% ============================================================
\section{Discretisation and Analysis}
% ============================================================

\subsection{Discrete Schemes}
[Derive your three methods step by step: Explicit Euler, RK4 (show the Butcher
tableau), and your implicit method with Newton. Number every equation.]

\begin{example}[ODE IVP: RK4 derivation]
For $\dot{y} = f(t,y)$, the classical Runge--Kutta method of order 4 is
\begin{align}
    k_1 &= f(t_n, y_n), \\
    k_2 &= f\!\left(t_n + \frac{h}{2}, y_n + \frac{h}{2}k_1\right), \\
    k_3 &= f\!\left(t_n + \frac{h}{2}, y_n + \frac{h}{2}k_2\right), \\
    k_4 &= f(t_n + h, y_n + hk_3), \\
    y_{n+1} &= y_n + \frac{h}{6}(k_1 + 2k_2 + 2k_3 + k_4).
\end{align}
\end{example}

\subsection{Consistency and Truncation Error}
[Show the Taylor expansion. State the order of accuracy of each method and under
what assumptions it holds.]

\subsection{Stability Analysis}
[The most critical section. Do not quote a textbook result --- apply the analysis
to \emph{your} problem. Derive the stability function $R(z)$ of each method, plot
the absolute-stability region, and connect it to the eigenvalues of your
linearised system.]

\begin{equation}
    R_{\text{Euler}}(z) = 1 + z, \qquad
    R_{\text{RK4}}(z) = 1 + z + \frac{z^2}{2} + \frac{z^3}{6} + \frac{z^4}{24},
    \qquad
    R_{\text{impl.Euler}}(z) = \frac{1}{1 - z}.
\end{equation}

\subsection{Convergence}
[Link consistency + stability to convergence (Lax equivalence or an energy
argument). State the theorem and verify your problem satisfies its hypotheses.]

\newpage
% ============================================================
\section{Implementation}
% ============================================================

\subsection{Algorithm}
[Present the core algorithm in pseudocode using the \texttt{algorithmic}
environment.]

\begin{algorithm}[H]
\caption{[Your method name]}
\label{alg:main}
\begin{algorithmic}[1]
\Require [initial data $y_0$, step size $h$, final time $T$]
\Ensure [approximate solution at $T$]
\State [Step 1]
\For{$n = 0, 1, \ldots, N-1$}
    \State [Loop body]
    \If{[condition]}
        \State [Action]
    \EndIf
\EndFor
\State \Return [result]
\end{algorithmic}
\end{algorithm}

\subsection{Code Structure}
[Briefly describe the organisation of your codebase. A table of files and their
purposes is helpful.]

\begin{table}[H]
\centering
\begin{tabular}{ll}
\toprule
\textbf{File} & \textbf{Purpose} \\
\midrule
\texttt{methods.py} & Core time-stepping: Euler, RK4, implicit Euler with Newton \\
\texttt{model.py} & The right-hand side $f(t,y)$ and its Jacobian \\
\texttt{adaptive.py} & Adaptive step-size controller \\
\texttt{run\_all.py} & Reproduces all report figures \\
\texttt{validation\_tests.py} & PASS/FAIL test suite \\
\bottomrule
\end{tabular}
\end{table}

\subsection{Key Design Decisions}
[Discuss non-obvious choices: why this implicit method, how Newton's initial
guess is chosen, how the adaptive controller estimates local error.]

\newpage
% ============================================================
\section{Numerical Experiments}
% ============================================================

\subsection{Verification: Method of Manufactured Solutions or Exact Solutions}
[Present a test case with a known solution. Show the error table and the
convergence plot.]

\begin{table}[H]
\centering
\begin{tabular}{cccc}
\toprule
\textbf{Step size $h$} & \textbf{Error} & \textbf{Ratio} & \textbf{Order} \\
\midrule
$0.1$    & [value] & ---  & ---  \\
$0.05$   & [value] & [ratio] & [order] \\
$0.025$  & [value] & [ratio] & [order] \\
\bottomrule
\end{tabular}
\caption{Convergence study for [test case].}
\end{table}

\subsection{Stability Experiments}
[Show that your explicit methods blow up beyond their stability limit while the
implicit method does not. Compare with the predicted threshold.]

\subsection{Adaptivity}
[Show the step size actually changes, the local error stays near the tolerance,
and the $h(t)$ trajectory.]

\subsection{Main Results and Comparison of Methods}
[Present the experiments that answer the project's central question. Use
high-quality figures with interpretive captions.]

\begin{figure}[H]
\centering
\fbox{\parbox{0.8\textwidth}{\centering \vspace{3em} [Insert your figure here] \vspace{3em}}}
\caption{[Caption that interprets, not merely describes. E.g., ``The log-log plot
confirms second-order convergence (slope 1.98) on the smooth test case, and shows
the round-off floor at $h \approx 10^{-6}$.'']}
\label{fig:convergence}
\end{figure}

\subsection{Limitations and Robustness}
[Discuss parameter regimes where your method struggles. This is not a weakness ---
it is scientific honesty.]

\newpage
% ============================================================
\section{Conclusions}
% ============================================================

[Summarise in 3--4 bullet points:]
\begin{itemize}[leftmargin=2em]
\item [What you achieved.]
\item [The most surprising or instructive finding.]
\item [The main limitation you identified.]
\item [A concrete direction for future work.]
\end{itemize}

\newpage
% ============================================================
\section*{References}
% ============================================================

[Use BibTeX or list references manually. Include at least: (1) the textbook from
which you learned the method; (2) a paper applying the method to a problem
similar to yours; (3) a software reference (NumPy, SciPy, scikit-fem).]

\begin{thebibliography}{9}
\bibitem{ref1} Author, A. \emph{Title}. Journal, year.
\bibitem{ref2} Author, B. \emph{Book title}. Publisher, year.
\end{thebibliography}

\newpage
% ============================================================
\appendix
\section{AI Transparency Log}
% ============================================================

[Append the completed team AI Transparency Log here. See
\texttt{06\_Templates/ICS} for the full form: tools used, significant claim audits,
what you changed and why, and the intellectual ownership declarations.]

\section{Individual Contribution Statements}
% ============================================================

Each team member must complete the following (repeat for all 5 members):

\begin{formbox}
\textbf{Name:} [Jingyan.Wei] \quad (\textbf{Student ID:} [2361411]) \hfill \textbf{Role:} Mathematical Theory (Schemes \& Stability Lead)

\vspace{0.5em}
\textbf{1. My specific contributions to this project:}

I authored the discrete formulation and linear/nonlinear stability foundations in Section~2 (Sections~2.1 and~2.3) in \texttt{Project1/report/Section2Draft.tex} (integrated into \texttt{report.tex}) and \texttt{Project1/notes/section2\_stability\_notes.md}. Specifically, I derived the step-update relations for Explicit Euler (Eq.~6--7 in Section2Draft\_v5), classical RK4 with its strictly lower-triangular Butcher tableau (Eq.~12--19), and Implicit Euler (Eq.~20--21). For the implicit solver, I derived the non-commutative matrix Fr{\'e}chet derivative $\mathcal{D}f(X)[H] = H(I_n - X^TX) - X(H^TX + X^TH)$ (Eq.~25) and the residual Jacobian $J_F(w) = I_{n^2} - h J_f(w)$ (Eq.~24, 28), formalising the damped Newton iteration with Armijo-type backtracking line search (Algorithm~1). Furthermore, I proved the analytical stability boundaries for the scalar Dahlquist equation: showing that $R_{\text{RK4}}(-a) = -1$ has no real roots (Eq.~76) and solving the cubic equation $q(a) = a^3 - 4a^2 + 12a - 24 = 0$ to establish the negative real intercept $a^* \approx 2.78529$ (Eq.~77) and imaginary limit $[-2\sqrt{2}, 2\sqrt{2}]$ (Eq.~78). For Topic~5, I transformed the linearised operator into SVD coordinates (Eq.~91), rigorously separating contracting normal modes $\lambda_{ij}^{\text{sym}}$ from neutral tangent modes $\lambda_{ij}^{\text{skew}}$ (Eq.~92--95), determined the full initial spectrum $\{-26.0, \dots, +0.88\}$ (Eq.~96), calculated the theoretical explicit step limits $h_{\text{crit}}^{\text{EE}} \approx 0.0769$ and $h_{\text{crit}}^{\text{RK4}} \approx 0.1071$ (Eq.~97--98), and authored the Mandatory Analytical Disclaimer governing local frozen spectra.

\vspace{0.5em}
\textbf{2. The hardest conceptual obstacle I encountered:}

The hardest mathematical challenge was deriving the problem-specific Fr{\'e}chet derivative and Jacobian spectrum for the matrix gradient flow. In standard scalar calculus, differentiating $x(1-x^2)$ yields $1-3x^2$, but for the matrix-valued flow $\dot{X} = X(I_n - X^TX)$, matrix multiplication is strictly non-commutative. When taking directional variations, the perturbation in the quadratic term expands as $\mathcal{D}(X^TX)[H] = H^TX + X^TH$ rather than $2X^TH$. Projecting this onto the tangent space of the SVD manifold required diagonalising a coupled $2 \times 2$ block system for every coordinate pair $(i < j)$, where I had to analytically demonstrate that symmetric perturbations induce contraction ($\lambda = -2$) while skew-symmetric variations align with the Lie algebra $\mathfrak{so}(n)$ of neutral tangent rotations ($\lambda = 0$).

\vspace{0.5em}
\textbf{3. One piece of AI-generated output my team used, and what I changed:}

Cross-referenced to \textbf{Appendix A (AI Transparency Log)}: Appendix~A documents that generative AI was used to detail the RK4 negative-real-axis stability proof and analyze numerical stability. When querying AI to evaluate explicit stability on Topic~5, the raw output claimed that ensuring the scaled frozen eigenvalues satisfied $|1 + h\lambda_j| \le 1$ ($h < 2/26 \approx 0.0769$) provided a sufficient mathematical proof of global numerical stability over $t \in [0, 2]$. I audited this claim and identified that because the Jacobian $J_f(X(t))$ drifts continuously along the nonlinear trajectory, frozen eigenvalues only provide a necessary local linearised indicator rather than a global nonlinear certificate. I revised Section~2.3.4 to introduce the Mandatory Analytical Disclaimer and specified the step-sweep perturbation protocol to empirically test true nonlinear blow-up boundaries.

\vspace{0.5em}
\textbf{4. One thing I still do not fully understand:}

I do not yet fully understand how the non-normality of the full $n^2 \times n^2$ vectorised Jacobian $J_f(X)$ affects transient perturbation growth before asymptotic decay sets in. While our spectrum analysis accurately tracks the eigenvalues of the symmetric transformed operator $\mathcal{L}_\Sigma$, the physical Kronecker-product Jacobian in column-major coordinates is generally non-normal ($J_f J_f^T \ne J_f^T J_f$), meaning that pseudospectral effects could theoretically induce transient growth even when all eigenvalues have negative real parts.

\vspace{0.8em}
\textbf{Signature:} \underline{\hspace{4.5cm}} \qquad \textbf{Date:} \underline{\hspace{2.5cm}}
\end{formbox}

\section{Supplementary Figures and Data}
% ============================================================

[Place any additional figures, tables, or code listings that would interrupt the
main narrative but may be useful to a curious reader.]

\end{document}
